\chapter{Metodologia}
\label{cap:metodologia}

A pesquisa combina a análise documental e a comparação técnica dos requisitos de conexão coletiva com a avaliação de seus possíveis impactos sobre o projeto elétrico aprovado do Edifício Palatium. A unidade de análise normativa corresponde às exigências relativas a empreendimentos com múltiplas unidades consumidoras na área de concessão da Enel Distribuição Ceará. O estudo de caso permite examinar as implicações das mudanças nos requisitos normativos sobre soluções de engenharia documentadas, sem constituir avaliação de conformidade integral do empreendimento.

\section{Natureza e abordagem da pesquisa}
\label{sec:tipo-pesquisa}

A pesquisa é classificada quanto à natureza, aos objetivos, à abordagem e aos procedimentos, conforme \citeonline[p.~31--39]{silveiracordova2009}. Quanto à natureza, trata-se de pesquisa aplicada, destinada a subsidiar decisões de projeto relacionadas à atualização das especificações de conexão coletiva. Quanto aos objetivos, caracteriza-se como descritiva, pois identifica e sistematiza as diferenças entre os requisitos das especificações 0126 e 1580 e examina suas possíveis repercussões técnicas.

A abordagem é predominantemente qualitativa, uma vez que a classificação das diferenças depende da interpretação do conteúdo das prescrições e de suas condições de aplicação. A análise é complementada por dados e cálculos quantitativos quando a avaliação dos impactos exige a comparação de limites de potência, demanda, corrente, dimensões ou características nominais dos equipamentos.

Quanto aos procedimentos, a pesquisa combina análise documental e estudo de caso único. A análise documental permite examinar as especificações técnicas e a documentação do projeto como fontes primárias, enquanto a comparação constitui o procedimento analítico utilizado para identificar as diferenças entre as prescrições \cite[p.~37]{silveiracordova2009}. A segunda etapa consiste na avaliação dessas diferenças em relação ao projeto elétrico do empreendimento. A delimitação do caso permite aprofundar a análise de um objeto específico, em conformidade com a discussão de \citeonline[p.~219 e 222]{miguel2007} sobre o estudo de caso em engenharia. Os resultados relativos ao empreendimento não são tomados como representativos de todos os projetos de conexão coletiva, embora a comparação normativa abranja os requisitos aplicáveis à área de concessão do Ceará.

A ausência de determinado requisito em uma seção não será interpretada, isoladamente, como sua eliminação, sendo necessária a verificação de remissões a outras seções, tabelas, anexos e documentos complementares.

\section{Conjunto documental}
\label{sec:documentos}

As fontes normativas principais são a especificação 0126, versão 03, de 31 de março de 2022, adotada como referência anterior do estudo de caso, e a especificação 1580, versão 2, de 8 de setembro de 2025, adotada como documento de comparação. A seleção fundamenta-se na indicação expressa de substituição da 0126 pelo novo documento corporativo. Conforme seu histórico de versões, a emissão inicial da 1580 ocorreu em 18 de agosto de 2025 (versão 1), e a versão 2 registra ajustes de formatação \cite[item~2, p.~3--4]{enel1580}. Adota-se a versão 2 por corresponder à revisão disponibilizada no catálogo de normas técnicas da Enel Distribuição Ceará na data da consulta \cite{enelnormasce}. Os documentos complementares serão consultados quando necessários à interpretação dos requisitos.

O histórico da 1580 menciona a transição normativa com remissão ao art.~19, §~2º, da Resolução Normativa ANEEL nº~1.000/2021 \cite[item~2, p.~4]{enel1580}. Entretanto, é o art.~20, §~2º, dessa resolução que estabelece, como regra geral, o prazo mínimo de 120 dias entre a comunicação das normas técnicas aos interessados e sua entrada em vigor, ressalvadas as hipóteses excepcionais previstas no próprio dispositivo \cite[art.~20, §~2º]{aneel1000}.

Considerando que a versão 1 da especificação 1580 foi emitida em 18 de agosto de 2025 e que a versão 2 apresentou apenas ajustes de formatação, estima-se que sua entrada em vigor tenha ocorrido em 16 de dezembro de 2025, correspondente ao término do prazo de 120 dias contado a partir da emissão inicial. Essa estimativa pressupõe que a comunicação oficial tenha coincidido com a emissão da primeira versão e que a distribuidora tenha adotado o prazo mínimo regulamentar. No entanto, a própria especificação 1580 não informa uma data específica de entrada em vigor, e o catálogo da distribuidora apresenta apenas a data de revisão de 8 de setembro de 2025 \cite[item~2, p.~3--4]{enel1580}\cite{enelnormasce}. Assim, adota-se 16 de dezembro de 2025 como referência temporal para a pesquisa, sem tratá-la como data de vigência oficialmente comprovada.

A Resolução Normativa ANEEL nº 1.000/2021 e suas alterações pertinentes compõem o referencial regulatório para interpretação das condições de atendimento, das normas técnicas da distribuidora e dos períodos de transição. Também integram o conjunto documental os desenhos e anexos das especificações quando contiverem exigências que complementem o texto principal.

O estudo de caso utiliza quatro documentos primários do empreendimento: o \emph{Memorial de cálculo elétrico dos apartamentos e condomínio}, composto por 13 páginas; o conjunto de cinco pranchas do projeto elétrico executivo, relacionadas na página~4 do memorial; a Anotação de Responsabilidade Técnica (ART) de obra/serviço, registrada em 7 de outubro de 2025, relativa ao projeto de instalações elétricas em baixa tensão e ao gerador de 81~kVA; e a carta da Enel Distribuição Ceará, com o assunto \emph{Carta resposta de análise em atenção a sua solicitação}, datada de 25 de outubro de 2025 \cite{palatiummemorial2025,palatiumpranchas2025,palatiumart2025,enelaprovacao2025}. O memorial e as pranchas apresentam carimbos de aprovação datados de 25 de outubro de 2025. A documentação aprovada constitui a base técnica da comparação; a ART identifica a atividade de elaboração do projeto, não constituindo, por si só, comprovação de responsabilidade técnica pela execução das instalações \cite[p.~1, item~4]{palatiumart2025}, e a carta formaliza a aprovação do projeto.

As pranchas 001 e 002 apresentam as plantas de alimentação do térreo e do subsolo, respectivamente; a prancha 003 apresenta o diagrama unifilar geral do centro de proteção geral, do centro de medição e dos apartamentos; a prancha 004 apresenta o diagrama unifilar do condomínio e do gerador; e a prancha 005 apresenta o esquema vertical de alimentação \cite[p.~4]{palatiummemorial2025}\cite[pranchas~001--005]{palatiumpranchas2025}. A identificação das evidências utiliza a página do memorial ou o número da prancha e o detalhe correspondente, conforme a natureza do dado analisado.

A identificação pública do estudo de caso restringe-se à denominação Edifício Palatium, ao município de Sobral e às características técnicas pertinentes à pesquisa. A autoria dos documentos particulares de projeto é identificada nos originais, mas não divulgada nas referências deste trabalho, que os individualizam por seus títulos, ano, local e extensão. O histórico interno de revisões não integra o objeto da pesquisa. A rastreabilidade dos dados é assegurada pela indicação das páginas, pranchas e elementos técnicos consultados. Não são reproduzidos os documentos integrais, nem incluídos identificadores pessoais, assinaturas, números de registros profissionais ou elementos de identificação da empresa projetista. A divulgação de eventuais recortes de documentos particulares depende da verificação das condições de utilização e reprodução aplicáveis, restringindo-se ao conteúdo indispensável à análise.

\section{Caracterização do empreendimento}
\label{sec:caracterizacao-empreendimento}

O Edifício Palatium é um empreendimento residencial multifamiliar vertical localizado em Sobral, Ceará. O projeto original contempla 26 apartamentos de 163~\(\mathrm{m}^2\), distribuídos em 13 pavimentos-tipo, com duas unidades por pavimento. O centro de medição único reúne 28 medições: 26 destinadas aos apartamentos, uma ao quadro de distribuição do condomínio (QDC) e uma ao quadro do sistema de bombeamento de incêndio \cite[p.~1 e~5]{palatiummemorial2025}.

A alimentação prevista consiste em uma entrada coletiva em baixa tensão trifásica de 380/220~V, derivada de rede de distribuição aérea, conforme informação técnica do autor, com ramal de entrada subterrâneo entre a entrada do empreendimento e o sistema de proteção e medição geral. O memorial e a prancha~001 documentam o encaminhamento subterrâneo dos condutores \cite[p.~1 e~7]{palatiummemorial2025}\cite[prancha~001]{palatiumpranchas2025}. O memorial declara, para o centro de proteção geral (CPG), carga instalada total de 707,31~kW e demanda de 291,33~kVA, com proteção geral por disjuntor termomagnético de 500~A. O projeto inclui um gerador de emergência de 81~kVA, trifásico, sem operação em paralelo com a rede da distribuidora \cite[p.~1 e~7]{palatiummemorial2025}\cite[pranchas~003 e~004]{palatiumpranchas2025}. Esses parâmetros descrevem a configuração projetada, não o estado das instalações executadas.

A carta da Enel Distribuição Ceará e os carimbos de aprovação do memorial e das pranchas comprovam a aprovação do projeto em 25 de outubro de 2025 \cite{enelaprovacao2025}\cite[p.~1]{palatiummemorial2025}\cite[pranchas~001--005]{palatiumpranchas2025}. Adota-se como premissa do estudo de caso que a aprovação ocorreu com base na especificação 0126, conforme as informações relativas ao projeto, uma vez que a carta e os carimbos não identificam expressamente o código da norma adotada. A escolha do empreendimento decorre da disponibilidade da documentação aprovada e da existência de soluções de conexão coletiva que permitem examinar, de forma documentada, as consequências técnicas da aplicação hipotética da 1580. Essa caracterização constitui a base de referência da análise; a comparação dos requisitos, os cálculos demonstrativos e a avaliação dos impactos são apresentados no Capítulo~\ref{cap:resultados}.

\section{Procedimentos de comparação normativa}
\label{sec:procedimento}

O levantamento será realizado em etapas sucessivas. Inicialmente, serão identificados objeto, área de aplicação e elementos técnicos equivalentes nas duas especificações. Em seguida, as prescrições serão agrupadas conforme os temas: enquadramento e limites de atendimento, cálculo de demanda, formas de conexão, ramais, padrão de entrada, medição, proteção, aterramento e exigências de apresentação de projeto.

O levantamento abrangerá os requisitos técnicos e documentais pertinentes à conexão coletiva na área de concessão do Ceará, independentemente de sua incidência sobre o empreendimento. Serão incluídas as prescrições gerais aplicáveis ao Ceará, as disposições específicas dessa área de concessão e as condições de atendimento pertinentes, inclusive quando um requisito estiver presente em apenas uma das especificações. Requisitos exclusivos de outras áreas de concessão serão excluídos, com registro do fundamento da exclusão. Quando a aplicabilidade territorial não puder ser estabelecida com segurança, a dúvida será registrada separadamente, sem presumir sua aplicação ao Ceará.

Para cada tema, serão registrados os dispositivos correspondentes em cada documento, com indicação da seção, item, tabela ou desenho, quando aplicável. A comparação será efetuada entre prescrições funcionalmente equivalentes, inclusive quando inseridas em partes distintas das especificações. Quando a 1580 remeter a outra norma, serão examinadas a referência e suas condições de aplicação antes de caracterizar eventual supressão ou modificação do requisito.

A unidade de classificação será o requisito técnico ou documental, considerado em suas condições de aplicação. Quando um mesmo dispositivo contiver prescrições independentes, estas serão analisadas separadamente, de modo que cada comparação receba uma única categoria. O Quadro~\ref{qua:classificacao-normativa} estabelece as definições operacionais e apresenta exemplos hipotéticos, que não constituem resultados da comparação das especificações.

\begin{quadro}[htbp]
  \centering
  \begin{Spacing}{1.}
    \caption{Classificação das alterações normativas}
    \label{qua:classificacao-normativa}
  \end{Spacing}
  {\small
  \renewcommand{\arraystretch}{1.2}
  \setlength{\tabcolsep}{4pt}
  \begin{tabular}{|>{\raggedright\arraybackslash}p{0.17\linewidth}|>{\raggedright\arraybackslash}p{0.43\linewidth}|>{\raggedright\arraybackslash}p{0.28\linewidth}|}
    \hline
    \textbf{Categoria} & \textbf{Critério de classificação} & \textbf{Exemplo hipotético} \\\hline
    Mantido & Prescrição equivalente quanto à exigência, aos parâmetros e às condições de aplicação, sem reorganização relevante. & Mesmo limite de corrente e mesmas condições de aplicação. \\\hline
    Alterado & Modificação de exigência, parâmetro, condição de aplicação ou solução admitida, inclusive substituição por requisito funcionalmente correspondente. & Alteração da seção mínima exigida de um condutor. \\\hline
    Acrescentado & Prescrição identificada na especificação 1580 sem equivalente na especificação 0126, após conferência dos documentos. & Inclusão da exigência de um documento de projeto antes não previsto na especificação. \\\hline
    Retirado & Prescrição da especificação 0126 cuja supressão na especificação 1580 é demonstrável, sem equivalente ou remissão que a preserve. & Supressão comprovada da exigência de um acessório específico. \\\hline
    Reorganizado & Mudança de localização, estrutura ou redação relevante, sem alteração técnica demonstrada, preservada a equivalência da prescrição. & Requisito transferido do texto para uma tabela, sem mudança de condições ou valores. \\\hline
    Inconclusivo & Evidências insuficientes para estabelecer equivalência, mudança, inclusão ou supressão com segurança. & Remissão genérica a documento não identificado, que não permite confirmar a equivalência. \\\hline
  \end{tabular}}
  \Fonte{Elaborado pelo autor.}
\end{quadro}

A classificação seguirá a seguinte lógica: inicialmente, serão confrontados os requisitos e suas condições de aplicação, considerando o texto, as notas, as tabelas, os desenhos, os anexos e as remissões identificáveis. Quando houver correspondência e for constatada mudança técnica, será adotada a categoria alterado. Se a equivalência for demonstrada, será adotada a categoria reorganizado quando houver mudança relevante na apresentação da prescrição, ou mantido nos demais casos. Na ausência de correspondente, a categoria acrescentado ou retirado dependerá da demonstração da inclusão ou da supressão, respectivamente. Se a evidência não permitir a distinção, será adotada a categoria inconclusivo; a simples ausência de um requisito em determinada seção não será suficiente para classificá-lo como retirado.

A Figura~\ref{fig:fluxograma-classificacao-normativa} sintetiza essa sequência decisória e complementa as definições operacionais do Quadro~\ref{qua:classificacao-normativa}. A verificação dos dispositivos e das remissões precederá a classificação, de modo que cada requisito receba uma única categoria e que a ausência em uma seção, isoladamente, não seja tomada como evidência de supressão.

A classificação refere-se exclusivamente às diferenças entre os requisitos das especificações 0126 e 1580, e não à data de constituição de obrigações no ordenamento regulatório. Uma prescrição incluída na 1580, sem equivalente na 0126 após a verificação dos dispositivos e das remissões aplicáveis, será classificada como acrescentado mesmo quando decorrer de obrigação regulatória anterior. Se o requisito já constar das duas especificações, será classificado conforme sua equivalência e organização, ou conforme eventual alteração das exigências e condições de aplicação. A existência de disposição regulatória anterior não determina, por si só, a categoria de classificação.

A justificativa de cada classificação será registrada juntamente com os dispositivos examinados, explicitando as diferenças observadas e as eventuais limitações documentais. Quando a origem regulatória anterior de um requisito for demonstrada, serão identificados o ato, o dispositivo e sua relação com as versões comparadas, sem criar categoria adicional. A avaliação dos impactos distinguirá uma exigência técnica efetivamente introduzida ou modificada pela 1580 da explicitação de uma obrigação que já incidia sobre o projeto. A existência de nota regional, condição excepcional ou exigência de estudo de viabilidade será registrada expressamente, de modo a evitar generalizações que não se apliquem ao Ceará. As categorias descrevem as diferenças entre as especificações e não indicam, por si sós, a necessidade de adequação do projeto utilizado no estudo de caso.

\begin{figure}[p]
  \centering
  \begin{Spacing}{1.}
    \caption{Procedimento de classificação das alterações normativas}
    \label{fig:fluxograma-classificacao-normativa}
  \end{Spacing}
  \noindent% Incluído pela seção de procedimentos de comparação normativa; requer tikz e as bibliotecas declaradas no preâmbulo.
% As coordenadas estão em centímetros. As seis categorias são mutuamente
% exclusivas; os dois caminhos inconclusivos convergem para a mesma saída.
\begin{tikzpicture}[
  x=1cm,y=1cm,
  font=\fontsize{9.5}{11}\selectfont,
  inicio/.style={draw,rounded corners=3mm,align=center,text width=8.8cm,
    minimum height=1.2cm,inner sep=2mm,line width=0.6pt},
  decisao/.style={draw,diamond,aspect=1.8,align=center,text width=3.35cm,
    inner sep=1mm,line width=0.6pt},
  resultado/.style={draw,rounded corners=2mm,align=center,text width=2.65cm,
    minimum height=0.85cm,inner sep=1.5mm,line width=0.6pt,font=\bfseries\fontsize{10}{12}\selectfont},
  seta/.style={-{Latex[length=2mm,width=1.4mm]},line width=0.6pt},
  rotulo/.style={font=\fontsize{9}{11}\selectfont,fill=white,inner sep=1.2pt,align=center}
]
  \node[inicio] (inicio) at (0,0) {\textbf{Início}\\
    Identificar e comparar os requisitos\\
    das especificações 0126 e 1580};
  \node[decisao] (evidencia) at (0,-2.5)
    {Há evidências documentais suficientes?};
  \node[resultado] (inconclusivo) at (5.3,-2.5) {Inconclusivo};
  \node[decisao] (correspondencia) at (0,-6.5)
    {Existe requisito funcionalmente correspondente na outra especificação?};

  \node[decisao] (tecnica) at (-4.1,-9.1)
    {Houve alteração técnica do requisito ou de suas condições de aplicação?};
  \node[decisao] (demonstracao) at (4.1,-9.1)
    {A inclusão ou supressão do requisito foi demonstrada?};
  \node[resultado] (alterado) at (-5.9,-11.3) {Alterado};

  \node[decisao] (reorganizacao) at (-4.1,-13.3)
    {Houve reorganização relevante da apresentação do requisito?};
  \node[decisao] (presenca) at (4.1,-13.3)
    {Em qual especificação está presente o requisito sem correspondente?};
  \node[resultado] (reorganizado) at (-5.9,-16.8) {Reorganizado};
  \node[resultado] (mantido) at (-2,-16.8) {Mantido};
  \node[resultado] (retirado) at (2,-16.8) {Retirado};
  \node[resultado] (acrescentado) at (5.9,-16.8) {Acrescentado};

  \draw[seta] (inicio.south) -- (evidencia.north);
  \draw[seta] (evidencia.east) -- node[rotulo,above] {Não} (inconclusivo.west);
  \draw[seta] (evidencia.south) -- node[rotulo,right] {Sim} (correspondencia.north);

  \draw[seta] (correspondencia.west) -| node[rotulo,pos=0.2,above] {Sim} (tecnica.north);
  \draw[seta] (correspondencia.east) -| node[rotulo,pos=0.2,above] {Não} (demonstracao.north);

  \draw[seta] (tecnica.west) -- ++(-0.6,0)
    |- node[rotulo,pos=0.25,left] {Sim} (alterado.west);
  \draw[seta] (tecnica.south) -- node[rotulo,right] {Não} (reorganizacao.north);

  % Rota externa: não atravessa os demais nós nem os rótulos das decisões.
  \draw[seta] (demonstracao.east) -- (7.6,-9.1)
    node[rotulo,midway,above] {Não}
    -- (7.6,-2.5) -- (inconclusivo.east);
  \draw[seta] (demonstracao.south) -- node[rotulo,right] {Sim} (presenca.north);

  \draw[seta] (reorganizacao.south west) -- ++(0,-0.45)
    -| node[rotulo,pos=0.6,left] {Sim} (reorganizado.north);
  \draw[seta] (reorganizacao.south east) -- ++(0,-0.45)
    -| node[rotulo,pos=0.6,right] {Não} (mantido.north);
  \draw[seta] (presenca.south west) -- ++(0,-0.45)
    -| node[rotulo,pos=0.6,left] {Apenas na\\especificação\\0126} (retirado.north);
  \draw[seta] (presenca.south east) -- ++(0,-0.45)
    -| node[rotulo,pos=0.6,right] {Apenas na\\especificação\\1580} (acrescentado.north);

  \node[align=center,text width=14.5cm,font=\fontsize{9}{11}\selectfont]
    at (0,-18) {Conferir texto, notas, tabelas, desenhos, anexos e remissões identificáveis.\\
    A ausência isolada em uma seção não comprova retirada.};
\end{tikzpicture}
\par
  \Fonte{Elaborado pelo autor.}
\end{figure}
% Mantém o fluxograma integralmente em uma página e dentro da seção de procedimentos de comparação normativa.
\clearpage

\section{Critérios de avaliação dos impactos técnicos}
\label{sec:impactos}

Os impactos serão avaliados segundo a possibilidade de interferência em decisões de engenharia documentadas. Serão considerados a modalidade de alimentação, a localização e a configuração do padrão de entrada, o dimensionamento de condutores e dispositivos de proteção, a disposição dos centros de medição e os requisitos de documentação e aprovação.

A seleção dos requisitos para o estudo de caso constituirá uma segunda etapa, posterior à comparação normativa. Para inclusão na avaliação dos impactos, deverão ser atendidos cumulativamente os seguintes critérios:

\begin{enumerate}
  \item aplicabilidade do requisito à área de concessão do Ceará;
  \item identificação de mudança de conteúdo técnico ou documental entre as especificações, inclusive alteração, inclusão ou supressão demonstrada;
  \item correspondência entre a mudança e uma característica, componente ou solução identificável na documentação do projeto do empreendimento.
\end{enumerate}

Os requisitos mantidos ou apenas reorganizados, sem alteração técnica demonstrada, não serão selecionados para essa avaliação. As diferenças que não incidirem sobre o empreendimento permanecerão registradas na comparação normativa geral, com a justificativa de sua não seleção. Quando a mudança normativa ou sua incidência sobre o empreendimento não puder ser estabelecida, a limitação será registrada na etapa correspondente, sem presumir ausência de impacto.

A seleção de um requisito não implica, por si só, necessidade de reformulação do projeto. Para cada requisito selecionado, serão confrontadas a solução documentada, as condições de aplicação e a exigência da especificação 1580 ou, no caso de supressão demonstrada, as consequências da ausência da exigência anterior. O Quadro~\ref{qua:classificacao-impactos} estabelece os critérios para a classificação dos resultados dessa avaliação.

\begin{quadro}[htbp]
  \centering
  \begin{Spacing}{1.}
    \caption{Classificação dos impactos técnicos no estudo de caso}
    \label{qua:classificacao-impactos}
  \end{Spacing}
  {\small
  \renewcommand{\arraystretch}{1.2}
  \setlength{\tabcolsep}{4pt}
  \begin{tabular}{|>{\raggedright\arraybackslash}p{0.29\linewidth}|>{\raggedright\arraybackslash}p{0.63\linewidth}|}
    \hline
    \textbf{Categoria de impacto} & \textbf{Critério de classificação} \\\hline
    Solução compatível & A solução documentada atende à exigência examinada, nas condições consideradas, ou a supressão demonstrada não requer modificação da solução para o atendimento à especificação 1580. \\\hline
    Adequação potencialmente necessária & A mudança normativa indica necessidade de modificar uma solução, parâmetro de dimensionamento, componente ou documento de projeto para atendimento à especificação 1580, nas condições analisadas. \\\hline
    Situação condicionada à avaliação da distribuidora & A conclusão depende de avaliação, estudo de viabilidade ou decisão da distribuidora expressamente prevista no requisito, não sendo possível determinar o atendimento apenas pela documentação disponível. \\\hline
    Informação documental insuficiente & O requisito e sua pertinência ao caso estão estabelecidos, mas faltam dados do projeto ou parâmetros indispensáveis para determinar a consequência técnica. \\\hline
  \end{tabular}}
  \Fonte{Elaborado pelo autor.}
\end{quadro}

A categoria informação documental insuficiente distingue-se da classificação normativa inconclusivo: nesta última, as evidências não permitem determinar a natureza da diferença entre as especificações; naquela, a mudança normativa e sua pertinência ao empreendimento foram identificadas, mas a consequência técnica não pode ser determinada com os dados disponíveis. Quando a insuficiência documental impedir a própria seleção do requisito, será registrada como limitação da seleção, sem atribuição de categoria de impacto.

Cada consequência será classificada separadamente. Quando for possível concluir a compatibilidade ou identificar uma adequação potencialmente necessária, será adotada a categoria correspondente. Caso a conclusão dependa de decisão expressamente reservada à distribuidora, será registrada a condição de avaliação, inclusive eventuais dados faltantes. Nos demais casos em que os dados não permitirem conclusão, será adotada a categoria informação documental insuficiente.

Quando disponíveis os parâmetros necessários, serão realizados cálculos demonstrativos com identificação das grandezas, hipóteses e fontes empregadas. Esses cálculos terão finalidade analítica e não substituirão os memoriais nem a responsabilidade técnica dos profissionais responsáveis pelo projeto original.

\section{Aplicação ao estudo de caso}
\label{sec:estudocaso}

A configuração do projeto original, apresentada na Seção~\ref{sec:caracterizacao-empreendimento}, constitui a referência para a avaliação hipotética das exigências selecionadas da especificação 1580.

A avaliação ficará restrita aos elementos presentes no memorial e nas pranchas do projeto. Para cada requisito selecionado segundo os critérios da Seção~\ref{sec:impactos}, será estabelecida uma relação entre a solução projetada, a evidência documental correspondente, a mudança normativa e a consequência técnica potencial. Serão registradas a justificativa da seleção, as hipóteses adotadas, a categoria atribuída e as limitações da avaliação. A pesquisa não promoverá alterações nos documentos originais nem constituirá laudo de conformidade ou reavaliação da aprovação do projeto.

A análise dos impactos não pressupõe a incidência retroativa da especificação 1580 sobre o projeto aprovado, nem se destina a estabelecer qual norma deveria ter sido adotada pela distribuidora na aprovação. Os resultados descrevem as consequências técnicas de sua aplicação hipotética à configuração documentada, sem depender da confirmação da data efetiva de transição.

\section{Tratamento dos resultados e limitações}
\label{sec:limitacoes}

A sistematização dos resultados baseia-se em uma matriz comparativa organizada por tema, na qual cada requisito é identificado por um código único. O Quadro~\ref{qua:campos-matriz} apresenta a estrutura das informações utilizadas na comparação normativa e na avaliação dos impactos técnicos.

\begin{quadro}[htbp]
  \centering
  \begin{Spacing}{1.}
    \caption{Estrutura da matriz comparativa}
    \label{qua:campos-matriz}
  \end{Spacing}
  {\small
  \renewcommand{\arraystretch}{1.2}
  \setlength{\tabcolsep}{4pt}
  \begin{tabular}{|>{\raggedright\arraybackslash}p{0.29\linewidth}|>{\raggedright\arraybackslash}p{0.63\linewidth}|}
    \hline
    \textbf{Grupo de campos} & \textbf{Informações registradas} \\\hline
    Identificação & Código do requisito e tema. \\\hline
    Especificação 0126 & Dispositivo de referência e descrição do requisito, incluindo condições de aplicação. \\\hline
    Especificação 1580 & Dispositivo correspondente e descrição do requisito, incluindo condições de aplicação. \\\hline
    Comparação normativa & Categoria definida no Quadro~\ref{qua:classificacao-normativa}, aplicabilidade ao Ceará e justificativa da classificação ou da exclusão. \\\hline
    Seleção para o estudo de caso & Decisão de seleção e justificativa, segundo os critérios da Seção~\ref{sec:impactos}; identificação dos registros de impacto associados. \\\hline
    Impacto no empreendimento & Solução documentada, consequência técnica, categoria de impacto definida no Quadro~\ref{qua:classificacao-impactos} e referência documental do projeto. \\\hline
    Sustentação da análise & Evidências normativas, hipóteses e incertezas, distintas das interpretações registradas na justificativa. \\\hline
  \end{tabular}}
  \Fonte{Elaborado pelo autor.}
\end{quadro}

Os requisitos são identificados pelos respectivos dispositivos das especificações, com indicação de seção, item, alínea, nota, tabela, desenho ou anexo, conforme a fonte. O registro distingue a descrição das prescrições e de suas condições de aplicação da interpretação que fundamenta a classificação. A ausência de requisito correspondente deve ser justificada com base nos dispositivos examinados.

Os impactos técnicos são vinculados ao código do requisito que lhes deu origem, permitindo avaliar diferentes consequências no empreendimento sem duplicar a classificação normativa. As diferenças aplicáveis ao Ceará que não incidem sobre o empreendimento integram a comparação geral, com a respectiva justificativa. Os requisitos excluídos e aqueles de aplicabilidade territorial indeterminada são identificados separadamente.

A consistência da análise é verificada mediante uma segunda leitura dos documentos originais, com conferência dos dispositivos, da correspondência entre os requisitos, das classificações atribuídas e dos critérios de seleção para o estudo de caso. A avaliação dos impactos também considera a correspondência com as soluções documentadas e a fundamentação das consequências técnicas, incluindo as hipóteses de cálculo. A conferência busca evitar duplicidades e preservar a coerência entre os registros e a síntese dos resultados. Esse procedimento não pressupõe validação independente.

O Capítulo~\ref{cap:resultados} reúne os principais resultados da comparação normativa e da avaliação dos impactos técnicos no empreendimento. O Apêndice A destina-se à apresentação integral da matriz comparativa e dos impactos associados, identificados pelos códigos dos requisitos.

A pesquisa limita-se às versões documentais selecionadas e aos critérios aplicáveis à área de concessão do Ceará. Não serão investigados os requisitos exclusivos de outras distribuidoras, o histórico de revisões internas do projeto, as instalações efetivamente executadas, os procedimentos de vistoria ou a energização do empreendimento. A análise não constitui laudo de conformidade nem substitui a avaliação técnica da distribuidora em situações sujeitas a estudos de viabilidade. Conforme a distinção de responsabilidades apresentada na Seção~\ref{sec:regulamentacao}, o objeto da pesquisa é a avaliação dos possíveis impactos da especificação 1580 sobre o projeto aprovado, tendo como referência a especificação 0126; não se pretende auditar a atuação dos profissionais responsáveis pelo projeto ou pela execução, nem os procedimentos de análise e aprovação da distribuidora.

Quando não houver elementos suficientes para determinada conclusão, a limitação será registrada no resultado correspondente. Esse procedimento preserva a distinção entre as informações expressamente documentadas, a interpretação de requisitos técnicos e as hipóteses dependentes de confirmação adicional.
